% Options for packages loaded elsewhere
\PassOptionsToPackage{unicode}{hyperref}
\PassOptionsToPackage{hyphens}{url}
%
\documentclass[
]{article}
\usepackage{amsmath,amssymb}
\usepackage{iftex}
\ifPDFTeX
  \usepackage[T1]{fontenc}
  \usepackage[utf8]{inputenc}
  \usepackage{textcomp} % provide euro and other symbols
\else % if luatex or xetex
  \usepackage{unicode-math} % this also loads fontspec
  \defaultfontfeatures{Scale=MatchLowercase}
  \defaultfontfeatures[\rmfamily]{Ligatures=TeX,Scale=1}
\fi
\usepackage{lmodern}
\ifPDFTeX\else
  % xetex/luatex font selection
\fi
% Use upquote if available, for straight quotes in verbatim environments
\IfFileExists{upquote.sty}{\usepackage{upquote}}{}
\IfFileExists{microtype.sty}{% use microtype if available
  \usepackage[]{microtype}
  \UseMicrotypeSet[protrusion]{basicmath} % disable protrusion for tt fonts
}{}
\makeatletter
\@ifundefined{KOMAClassName}{% if non-KOMA class
  \IfFileExists{parskip.sty}{%
    \usepackage{parskip}
  }{% else
    \setlength{\parindent}{0pt}
    \setlength{\parskip}{6pt plus 2pt minus 1pt}}
}{% if KOMA class
  \KOMAoptions{parskip=half}}
\makeatother
\usepackage{xcolor}
\usepackage[margin=1in]{geometry}
\usepackage{color}
\usepackage{fancyvrb}
\newcommand{\VerbBar}{|}
\newcommand{\VERB}{\Verb[commandchars=\\\{\}]}
\DefineVerbatimEnvironment{Highlighting}{Verbatim}{commandchars=\\\{\}}
% Add ',fontsize=\small' for more characters per line
\usepackage{framed}
\definecolor{shadecolor}{RGB}{248,248,248}
\newenvironment{Shaded}{\begin{snugshade}}{\end{snugshade}}
\newcommand{\AlertTok}[1]{\textcolor[rgb]{0.94,0.16,0.16}{#1}}
\newcommand{\AnnotationTok}[1]{\textcolor[rgb]{0.56,0.35,0.01}{\textbf{\textit{#1}}}}
\newcommand{\AttributeTok}[1]{\textcolor[rgb]{0.13,0.29,0.53}{#1}}
\newcommand{\BaseNTok}[1]{\textcolor[rgb]{0.00,0.00,0.81}{#1}}
\newcommand{\BuiltInTok}[1]{#1}
\newcommand{\CharTok}[1]{\textcolor[rgb]{0.31,0.60,0.02}{#1}}
\newcommand{\CommentTok}[1]{\textcolor[rgb]{0.56,0.35,0.01}{\textit{#1}}}
\newcommand{\CommentVarTok}[1]{\textcolor[rgb]{0.56,0.35,0.01}{\textbf{\textit{#1}}}}
\newcommand{\ConstantTok}[1]{\textcolor[rgb]{0.56,0.35,0.01}{#1}}
\newcommand{\ControlFlowTok}[1]{\textcolor[rgb]{0.13,0.29,0.53}{\textbf{#1}}}
\newcommand{\DataTypeTok}[1]{\textcolor[rgb]{0.13,0.29,0.53}{#1}}
\newcommand{\DecValTok}[1]{\textcolor[rgb]{0.00,0.00,0.81}{#1}}
\newcommand{\DocumentationTok}[1]{\textcolor[rgb]{0.56,0.35,0.01}{\textbf{\textit{#1}}}}
\newcommand{\ErrorTok}[1]{\textcolor[rgb]{0.64,0.00,0.00}{\textbf{#1}}}
\newcommand{\ExtensionTok}[1]{#1}
\newcommand{\FloatTok}[1]{\textcolor[rgb]{0.00,0.00,0.81}{#1}}
\newcommand{\FunctionTok}[1]{\textcolor[rgb]{0.13,0.29,0.53}{\textbf{#1}}}
\newcommand{\ImportTok}[1]{#1}
\newcommand{\InformationTok}[1]{\textcolor[rgb]{0.56,0.35,0.01}{\textbf{\textit{#1}}}}
\newcommand{\KeywordTok}[1]{\textcolor[rgb]{0.13,0.29,0.53}{\textbf{#1}}}
\newcommand{\NormalTok}[1]{#1}
\newcommand{\OperatorTok}[1]{\textcolor[rgb]{0.81,0.36,0.00}{\textbf{#1}}}
\newcommand{\OtherTok}[1]{\textcolor[rgb]{0.56,0.35,0.01}{#1}}
\newcommand{\PreprocessorTok}[1]{\textcolor[rgb]{0.56,0.35,0.01}{\textit{#1}}}
\newcommand{\RegionMarkerTok}[1]{#1}
\newcommand{\SpecialCharTok}[1]{\textcolor[rgb]{0.81,0.36,0.00}{\textbf{#1}}}
\newcommand{\SpecialStringTok}[1]{\textcolor[rgb]{0.31,0.60,0.02}{#1}}
\newcommand{\StringTok}[1]{\textcolor[rgb]{0.31,0.60,0.02}{#1}}
\newcommand{\VariableTok}[1]{\textcolor[rgb]{0.00,0.00,0.00}{#1}}
\newcommand{\VerbatimStringTok}[1]{\textcolor[rgb]{0.31,0.60,0.02}{#1}}
\newcommand{\WarningTok}[1]{\textcolor[rgb]{0.56,0.35,0.01}{\textbf{\textit{#1}}}}
\usepackage{graphicx}
\makeatletter
\def\maxwidth{\ifdim\Gin@nat@width>\linewidth\linewidth\else\Gin@nat@width\fi}
\def\maxheight{\ifdim\Gin@nat@height>\textheight\textheight\else\Gin@nat@height\fi}
\makeatother
% Scale images if necessary, so that they will not overflow the page
% margins by default, and it is still possible to overwrite the defaults
% using explicit options in \includegraphics[width, height, ...]{}
\setkeys{Gin}{width=\maxwidth,height=\maxheight,keepaspectratio}
% Set default figure placement to htbp
\makeatletter
\def\fps@figure{htbp}
\makeatother
\setlength{\emergencystretch}{3em} % prevent overfull lines
\providecommand{\tightlist}{%
  \setlength{\itemsep}{0pt}\setlength{\parskip}{0pt}}
\setcounter{secnumdepth}{-\maxdimen} % remove section numbering
\usepackage{booktabs}
\usepackage{longtable}
\usepackage{array}
\usepackage{multirow}
\usepackage{wrapfig}
\usepackage{float}
\usepackage{colortbl}
\usepackage{pdflscape}
\usepackage{tabu}
\usepackage{threeparttable}
\usepackage{threeparttablex}
\usepackage[normalem]{ulem}
\usepackage{makecell}
\usepackage{xcolor}
\ifLuaTeX
  \usepackage{selnolig}  % disable illegal ligatures
\fi
\usepackage{bookmark}
\IfFileExists{xurl.sty}{\usepackage{xurl}}{} % add URL line breaks if available
\urlstyle{same}
\hypersetup{
  pdftitle={Lab 1 Forecasting with ARIMA and ETS models},
  pdfauthor={Team 3: Ben Mahklouf, Tina Tao, and Emma Christman},
  hidelinks,
  pdfcreator={LaTeX via pandoc}}

\title{Lab 1 Forecasting with ARIMA and ETS models}
\author{Team 3: Ben Mahklouf, Tina Tao, and Emma Christman}
\date{}

\begin{document}
\maketitle

{
\setcounter{tocdepth}{2}
\tableofcontents
}
\#Initial Plan and What Actually Happened We will compare the accuracy
of ARIMA models in yearly vs.~monthly data, answering the question: Is
the accuracy of monthly or yearly data better? To do this we will look
only at the Washington data. To begin, we will test the stationarity of
the yearly data and the monthly data.

If either are not stationary, we will use differencing to make them
stationary for the purposes of this analysis. After the data are
stationary, we will fit an ARIMA and ETS model to the monthly data using
AutoARIMA on the earliest 80\% of the data. This will allow use to
compare the accuracy of the ARIMA model in predicting the last 20\% of
the data to the actual data. After this is done, we will do the same to
the yearly data. Then, we can use the accuracy function to compare the
two methods.

We loaded the data and wrangled it, filtering it to only include the
Washington data. This allowed us to compare between both the yearly and
monthly data. From there we used the auto.arima() function, which
auto-checked for stationarity for us. We needed to make the yearly data
stationary by differencing it once. Then, we fit both ETS and arima
models to both the monthly and yearly catch data from Washington. After
fitting the models, we concluded that the monthly data forecast was more
accurate than the yearly data forecast.`

\#Question Addressed 1) Is an ARIMA model and/or an ETS model adequate
for using monthly data and yearly data? 2) Is the accuracy of
forecasting with monthly data or yearly data ``better''?

\section{Step1, load data and filter to Washington with basic
plots}\label{step1-load-data-and-filter-to-washington-with-basic-plots}

\begin{Shaded}
\begin{Highlighting}[]
\CommentTok{\# Seperate the monthly and yearly data }
\FunctionTok{library}\NormalTok{(atsalibrary) }\CommentTok{\# This is ALL the data }
\end{Highlighting}
\end{Shaded}

\begin{verbatim}
## Loading required package: broom
\end{verbatim}

\begin{verbatim}
## Loading required package: corrplot
\end{verbatim}

\begin{verbatim}
## corrplot 0.95 loaded
\end{verbatim}

\begin{verbatim}
## Loading required package: depmixS4
\end{verbatim}

\begin{verbatim}
## Loading required package: nnet
\end{verbatim}

\begin{verbatim}
## Loading required package: MASS
\end{verbatim}

\begin{verbatim}
## 
## Attaching package: 'MASS'
\end{verbatim}

\begin{verbatim}
## The following object is masked from 'package:dplyr':
## 
##     select
\end{verbatim}

\begin{verbatim}
## Loading required package: Rsolnp
\end{verbatim}

\begin{verbatim}
## Loading required package: nlme
\end{verbatim}

\begin{verbatim}
## 
## Attaching package: 'nlme'
\end{verbatim}

\begin{verbatim}
## The following object is masked from 'package:dplyr':
## 
##     collapse
\end{verbatim}

\begin{verbatim}
## Loading required package: forecast
\end{verbatim}

\begin{verbatim}
## Registered S3 method overwritten by 'quantmod':
##   method            from
##   as.zoo.data.frame zoo
\end{verbatim}

\begin{verbatim}
## 
## Attaching package: 'forecast'
\end{verbatim}

\begin{verbatim}
## The following object is masked from 'package:nlme':
## 
##     getResponse
\end{verbatim}

\begin{verbatim}
## Loading required package: fpp2
\end{verbatim}

\begin{verbatim}
## -- Attaching packages ---------------------------------------------- fpp2 2.5 --
\end{verbatim}

\begin{verbatim}
## v fma       2.5     v expsmooth 2.3
\end{verbatim}

\begin{verbatim}
## -- Conflicts ------------------------------------------------- fpp2_conflicts --
## x forecast::getResponse() masks nlme::getResponse()
\end{verbatim}

\begin{verbatim}
## Loading required package: ggmap
\end{verbatim}

\begin{verbatim}
## i Google's Terms of Service: <https://mapsplatform.google.com>
##   Stadia Maps' Terms of Service: <https://stadiamaps.com/terms-of-service/>
##   OpenStreetMap's Tile Usage Policy: <https://operations.osmfoundation.org/policies/tiles/>
## i Please cite ggmap if you use it! Use `citation("ggmap")` for details.
## Loading required package: grid
## 
## Loading required package: gridExtra
## 
## 
## Attaching package: 'gridExtra'
## 
## 
## The following object is masked from 'package:dplyr':
## 
##     combine
## 
## 
## Loading required package: Hmisc
## 
## 
## Attaching package: 'Hmisc'
## 
## 
## The following objects are masked from 'package:dplyr':
## 
##     src, summarize
## 
## 
## The following objects are masked from 'package:base':
## 
##     format.pval, units
## 
## 
## Loading required package: kableExtra
## 
## 
## Attaching package: 'kableExtra'
## 
## 
## The following object is masked from 'package:dplyr':
## 
##     group_rows
## 
## 
## Loading required package: knitr
## 
## Loading required package: loo
## 
## This is loo version 2.8.0
## 
## - Online documentation and vignettes at mc-stan.org/loo
## 
## - As of v2.0.0 loo defaults to 1 core but we recommend using as many as possible. Use the 'cores' argument or set options(mc.cores = NUM_CORES) for an entire session. 
## 
## 
## Attaching package: 'loo'
## 
## 
## The following object is masked from 'package:fma':
## 
##     milk
## 
## 
## Loading required package: magrittr
## 
## 
## Attaching package: 'magrittr'
## 
## 
## The following object is masked from 'package:ggmap':
## 
##     inset
## 
## 
## The following object is masked from 'package:purrr':
## 
##     set_names
## 
## 
## The following object is masked from 'package:tidyr':
## 
##     extract
## 
## 
## Loading required package: MARSS
## 
## Loading required package: quantmod
## 
## Loading required package: xts
## 
## Loading required package: zoo
## 
## 
## Attaching package: 'zoo'
## 
## 
## The following objects are masked from 'package:base':
## 
##     as.Date, as.Date.numeric
## 
## 
## 
## ######################### Warning from 'xts' package ##########################
## #                                                                             #
## # The dplyr lag() function breaks how base R's lag() function is supposed to  #
## # work, which breaks lag(my_xts). Calls to lag(my_xts) that you type or       #
## # source() into this session won't work correctly.                            #
## #                                                                             #
## # Use stats::lag() to make sure you're not using dplyr::lag(), or you can add #
## # conflictRules('dplyr', exclude = 'lag') to your .Rprofile to stop           #
## # dplyr from breaking base R's lag() function.                                #
## #                                                                             #
## # Code in packages is not affected. It's protected by R's namespace mechanism #
## # Set `options(xts.warn_dplyr_breaks_lag = FALSE)` to suppress this warning.  #
## #                                                                             #
## ###############################################################################
## 
## 
## Attaching package: 'xts'
## 
## 
## The following objects are masked from 'package:dplyr':
## 
##     first, last
## 
## 
## Loading required package: TTR
## 
## 
## Attaching package: 'quantmod'
## 
## 
## The following object is masked from 'package:Hmisc':
## 
##     Lag
## 
## 
## Loading required package: RCurl
## 
## 
## Attaching package: 'RCurl'
## 
## 
## The following object is masked from 'package:tidyr':
## 
##     complete
## 
## 
## Loading required package: reshape2
## 
## 
## Attaching package: 'reshape2'
## 
## 
## The following object is masked from 'package:tidyr':
## 
##     smiths
## 
## 
## Loading required package: rmarkdown
## 
## Loading required package: tseries
## 
## Loading required package: urca
\end{verbatim}

\begin{Shaded}
\begin{Highlighting}[]
\CommentTok{\# This is the monthly data }
\NormalTok{chinook.month }\OtherTok{\textless{}{-}}\NormalTok{ chinook.month }\SpecialCharTok{\%\textgreater{}\%}
  \FunctionTok{mutate}\NormalTok{(}\AttributeTok{Date =} \FunctionTok{as.Date}\NormalTok{(}\FunctionTok{paste}\NormalTok{(Year, Month, }\StringTok{"01"}\NormalTok{, }\AttributeTok{sep =} \StringTok{"{-}"}\NormalTok{), }\AttributeTok{format =} \StringTok{"\%Y{-}\%b{-}\%d"}\NormalTok{))}

\NormalTok{chinook.month }\OtherTok{\textless{}{-}}\NormalTok{ chinook.month }\SpecialCharTok{\%\textgreater{}\%}
  \FunctionTok{filter}\NormalTok{(State }\SpecialCharTok{\%in\%} \FunctionTok{c}\NormalTok{(}\StringTok{"WA"}\NormalTok{))}

\NormalTok{chinook.year }\OtherTok{\textless{}{-}}\NormalTok{ chinook.year }

\NormalTok{chinook.year\_filtered }\OtherTok{\textless{}{-}}\NormalTok{ chinook.year }\SpecialCharTok{\%\textgreater{}\%}
  \FunctionTok{filter}\NormalTok{(State }\SpecialCharTok{==} \StringTok{"Washington"}\NormalTok{)}

\CommentTok{\# Quick plot of year }

\FunctionTok{ggplot}\NormalTok{(chinook.year\_filtered, }\FunctionTok{aes}\NormalTok{(}\AttributeTok{x =}\NormalTok{ Year, }\AttributeTok{y =}\NormalTok{ log.metric.tons)) }\SpecialCharTok{+}
  \FunctionTok{geom\_line}\NormalTok{() }\SpecialCharTok{+}
  \FunctionTok{facet\_wrap}\NormalTok{(}\SpecialCharTok{\textasciitilde{}}\NormalTok{ State, }\AttributeTok{ncol =} \DecValTok{1}\NormalTok{) }\SpecialCharTok{+}  \CommentTok{\# 3 panels vertically}
  \FunctionTok{labs}\NormalTok{(}\AttributeTok{title =} \StringTok{"Monthly Chinook Salmon Landings"}\NormalTok{,}
       \AttributeTok{x =} \StringTok{"Date"}\NormalTok{,}
       \AttributeTok{y =} \StringTok{"Log(Metric Tons)"}\NormalTok{) }\SpecialCharTok{+}
  \FunctionTok{theme\_minimal}\NormalTok{()}
\end{Highlighting}
\end{Shaded}

\includegraphics{Lab1-ARIMA-2025-Team3_files/figure-latex/unnamed-chunk-2-1.pdf}

\begin{Shaded}
\begin{Highlighting}[]
\CommentTok{\# Quick plot of month }

\FunctionTok{ggplot}\NormalTok{(chinook.month, }\FunctionTok{aes}\NormalTok{(}\AttributeTok{x =}\NormalTok{ Date, }\AttributeTok{y =}\NormalTok{ log.metric.tons)) }\SpecialCharTok{+}
  \FunctionTok{geom\_line}\NormalTok{() }\SpecialCharTok{+}
  \FunctionTok{facet\_wrap}\NormalTok{(}\SpecialCharTok{\textasciitilde{}}\NormalTok{ State, }\AttributeTok{ncol =} \DecValTok{1}\NormalTok{) }\SpecialCharTok{+}  \CommentTok{\# 3 panels vertically}
  \FunctionTok{labs}\NormalTok{(}\AttributeTok{title =} \StringTok{"Monthly Chinook Salmon Landings"}\NormalTok{,}
       \AttributeTok{x =} \StringTok{"Date"}\NormalTok{,}
       \AttributeTok{y =} \StringTok{"Log(Metric Tons)"}\NormalTok{) }\SpecialCharTok{+}
  \FunctionTok{theme\_minimal}\NormalTok{()}
\end{Highlighting}
\end{Shaded}

\includegraphics{Lab1-ARIMA-2025-Team3_files/figure-latex/unnamed-chunk-2-2.pdf}

\section{Step2, check stationarity and
differencing}\label{step2-check-stationarity-and-differencing}

\begin{Shaded}
\begin{Highlighting}[]
\NormalTok{yearly }\OtherTok{\textless{}{-}}\NormalTok{ chinook.year\_filtered }\SpecialCharTok{\%\textgreater{}\%} 
  \FunctionTok{arrange}\NormalTok{(Year) }\SpecialCharTok{\%\textgreater{}\%}
  \FunctionTok{pull}\NormalTok{(log.metric.tons)}

\NormalTok{yearly\_ts }\OtherTok{\textless{}{-}} \FunctionTok{ts}\NormalTok{(yearly, }\AttributeTok{start =} \DecValTok{1950}\NormalTok{, }\AttributeTok{end =} \DecValTok{2016}\NormalTok{, }\AttributeTok{frequency =} \DecValTok{1}\NormalTok{)}


\NormalTok{monthly}\OtherTok{\textless{}{-}}\NormalTok{ chinook.month }\SpecialCharTok{\%\textgreater{}\%}
  \CommentTok{\#change log metric tons inf to NA}
  \FunctionTok{mutate}\NormalTok{(}\AttributeTok{log.metric.tons =} \FunctionTok{ifelse}\NormalTok{(log.metric.tons }\SpecialCharTok{==} \ConstantTok{Inf}\NormalTok{, }\ConstantTok{NA}\NormalTok{, log.metric.tons)) }\SpecialCharTok{\%\textgreater{}\%}
  \FunctionTok{mutate}\NormalTok{(}\AttributeTok{log.metric.tons =} \FunctionTok{ifelse}\NormalTok{(log.metric.tons }\SpecialCharTok{==} \SpecialCharTok{{-}}\ConstantTok{Inf}\NormalTok{, }\ConstantTok{NA}\NormalTok{, log.metric.tons)) }\SpecialCharTok{\%\textgreater{}\%}
  \FunctionTok{filter}\NormalTok{(}\SpecialCharTok{!}\FunctionTok{is.na}\NormalTok{(log.metric.tons))}

\NormalTok{monthly }\OtherTok{\textless{}{-}}\NormalTok{ monthly }\SpecialCharTok{\%\textgreater{}\%}
  \FunctionTok{arrange}\NormalTok{(Date) }\SpecialCharTok{\%\textgreater{}\%}
  \FunctionTok{pull}\NormalTok{(log.metric.tons)}

\CommentTok{\# Are there NAs in monthly }
\CommentTok{\# remove NA}

\NormalTok{monthly\_ts }\OtherTok{\textless{}{-}} \FunctionTok{ts}\NormalTok{(monthly, }\AttributeTok{start =} \FunctionTok{c}\NormalTok{(}\DecValTok{1990}\NormalTok{,}\DecValTok{1}\NormalTok{), }\AttributeTok{frequency =} \DecValTok{12}\NormalTok{)}


\DocumentationTok{\#\#\#\#\#\#\#\#\#\#\#\#\#\#\#\#\#\#\#\#\#\#\#\#\#\#\#\#\#\#}

\CommentTok{\# Check the need to difference the data }

\FunctionTok{adf.test}\NormalTok{(yearly\_ts)}
\end{Highlighting}
\end{Shaded}

\begin{verbatim}
## 
##  Augmented Dickey-Fuller Test
## 
## data:  yearly_ts
## Dickey-Fuller = -2.5627, Lag order = 4, p-value = 0.3467
## alternative hypothesis: stationary
\end{verbatim}

\begin{Shaded}
\begin{Highlighting}[]
\CommentTok{\# Augmented Dickey{-}Fuller Test}
\CommentTok{\# }
\CommentTok{\# data:  yearly\_ts}
\CommentTok{\# Dickey{-}Fuller = {-}2.5627, Lag order = 4, p{-}value = 0.3467}
\CommentTok{\# alternative hypothesis: stationary}

\CommentTok{\# Not stationary, p value greater than .01 }


\FunctionTok{adf.test}\NormalTok{(monthly\_ts)}
\end{Highlighting}
\end{Shaded}

\begin{verbatim}
## 
##  Augmented Dickey-Fuller Test
## 
## data:  monthly_ts
## Dickey-Fuller = -9.6198, Lag order = 7, p-value = 0.01
## alternative hypothesis: stationary
\end{verbatim}

\begin{Shaded}
\begin{Highlighting}[]
\CommentTok{\# Augmented Dickey{-}Fuller Test}
\CommentTok{\# }
\CommentTok{\# data:  monthly\_ts}
\CommentTok{\# Dickey{-}Fuller = {-}9.6198, Lag order = 7, p{-}value = 0.01}
\CommentTok{\# alternative hypothesis: stationary}

\CommentTok{\# These data ARE Stationary }

\DocumentationTok{\#\#\#\#\#\#\#\#\#\#\#\#\#\#\#\#\#\#\#\#\#\#\#\#\#\#\#\#\#}

\CommentTok{\# We can run the ndiff to see how many differences we would need to get to stationarity }

\FunctionTok{ndiffs}\NormalTok{(yearly\_ts)}
\end{Highlighting}
\end{Shaded}

\begin{verbatim}
## [1] 1
\end{verbatim}

\begin{Shaded}
\begin{Highlighting}[]
\CommentTok{\# [1], we need to difference ONCE, we expect the middle number in autoarima to be 1 }
\end{Highlighting}
\end{Shaded}

\subsubsection{Step 3, explore correlation and autocorrelation using ACF
and
PACF}\label{step-3-explore-correlation-and-autocorrelation-using-acf-and-pacf}

\begin{Shaded}
\begin{Highlighting}[]
\FunctionTok{acf}\NormalTok{(yearly\_ts, }\AttributeTok{lag.max =} \DecValTok{20}\NormalTok{, }\AttributeTok{main =} \StringTok{"ACF of Yearly Chinook Salmon Landings in Washington"}\NormalTok{)}
\end{Highlighting}
\end{Shaded}

\includegraphics{Lab1-ARIMA-2025-Team3_files/figure-latex/unnamed-chunk-4-1.pdf}

\begin{Shaded}
\begin{Highlighting}[]
\CommentTok{\# The ACF plot shows a tails off structure}
\FunctionTok{pacf}\NormalTok{(yearly\_ts, }\AttributeTok{lag.max =} \DecValTok{20}\NormalTok{, }\AttributeTok{main =} \StringTok{"PACF of Yearly Chinook Salmon Landings in Washington"}\NormalTok{)}
\end{Highlighting}
\end{Shaded}

\includegraphics{Lab1-ARIMA-2025-Team3_files/figure-latex/unnamed-chunk-4-2.pdf}

\begin{Shaded}
\begin{Highlighting}[]
\CommentTok{\# The PACF plot show a cuts off structure after lag=1}
\CommentTok{\# These are AR(1) model characteristics}
\CommentTok{\# Combining with results from the above stationary test, we expect an ARIMA(1,1,0) model to fit the yearly Chinook Salmon Landings in Washington}

\FunctionTok{acf}\NormalTok{(monthly\_ts, }\AttributeTok{lag.max =} \DecValTok{36}\NormalTok{, }\AttributeTok{main =} \StringTok{"ACF of Monthly Chinook Salmon Landings in Washington"}\NormalTok{)}
\end{Highlighting}
\end{Shaded}

\includegraphics{Lab1-ARIMA-2025-Team3_files/figure-latex/unnamed-chunk-4-3.pdf}

\begin{Shaded}
\begin{Highlighting}[]
\CommentTok{\# We observe a clear deterministic seasonal pattern and expect a big negative coefficient for AR component}
\FunctionTok{pacf}\NormalTok{(monthly\_ts, }\AttributeTok{lag.max =} \DecValTok{36}\NormalTok{, }\AttributeTok{main =} \StringTok{"PACF of Monthly Chinook Salmon Landings in Washington"}\NormalTok{)}
\end{Highlighting}
\end{Shaded}

\includegraphics{Lab1-ARIMA-2025-Team3_files/figure-latex/unnamed-chunk-4-4.pdf}

\begin{Shaded}
\begin{Highlighting}[]
\CommentTok{\# Several correlation patterns are detected and there is a strong correlation with lag=2}
\CommentTok{\# Choosing the optimal p and q value may need to be done by comparing AIC and BIC values, which is expected to be conducted in auto.arima function}
\end{Highlighting}
\end{Shaded}

\subsubsection{Step 4, fit ARIMA and ETS models using created training
and test sets on the monthly and yearly
data}\label{step-4-fit-arima-and-ets-models-using-created-training-and-test-sets-on-the-monthly-and-yearly-data}

\begin{Shaded}
\begin{Highlighting}[]
\CommentTok{\# Create training and test sets (80\% training, 20\% test)}
\NormalTok{yearly\_train\_end }\OtherTok{\textless{}{-}} \FunctionTok{floor}\NormalTok{(}\FloatTok{0.8} \SpecialCharTok{*} \FunctionTok{length}\NormalTok{(yearly\_ts))}
\NormalTok{yearly\_train }\OtherTok{\textless{}{-}} \FunctionTok{window}\NormalTok{(yearly\_ts, }\AttributeTok{end =} \FunctionTok{time}\NormalTok{(yearly\_ts)[yearly\_train\_end])}
\NormalTok{yearly\_test }\OtherTok{\textless{}{-}} \FunctionTok{window}\NormalTok{(yearly\_ts, }\AttributeTok{start =} \FunctionTok{time}\NormalTok{(yearly\_ts)[yearly\_train\_end }\SpecialCharTok{+} \DecValTok{1}\NormalTok{])}

\NormalTok{monthly\_train\_end }\OtherTok{\textless{}{-}} \FunctionTok{floor}\NormalTok{(}\FloatTok{0.8} \SpecialCharTok{*} \FunctionTok{length}\NormalTok{(monthly\_ts))}
\NormalTok{monthly\_train }\OtherTok{\textless{}{-}} \FunctionTok{window}\NormalTok{(monthly\_ts, }\AttributeTok{end =} \FunctionTok{time}\NormalTok{(monthly\_ts)[monthly\_train\_end])}
\NormalTok{monthly\_test }\OtherTok{\textless{}{-}} \FunctionTok{window}\NormalTok{(monthly\_ts, }\AttributeTok{start =} \FunctionTok{time}\NormalTok{(monthly\_ts)[monthly\_train\_end }\SpecialCharTok{+} \DecValTok{1}\NormalTok{])}

\CommentTok{\# Delete monthly\_train\_end}
\FunctionTok{rm}\NormalTok{(monthly\_train\_end)}
\CommentTok{\# Delete yearly\_train\_end}
\FunctionTok{rm}\NormalTok{(yearly\_train\_end)}

\CommentTok{\# Fit ARIMA model to training data}
\NormalTok{yearly\_arima }\OtherTok{\textless{}{-}} \FunctionTok{auto.arima}\NormalTok{(yearly\_train)}
\CommentTok{\# Series: yearly\_train }
\CommentTok{\# The model result suggests an random walk ARIMA(0,1,0) instead of ARIMA(1,1,0), because the AR(1) coefficient might be small or not statistically significant}
\CommentTok{\# }
\CommentTok{\# sigma\^{}2 = 0.04824:  log likelihood = 5.04}
\CommentTok{\# AIC={-}8.07   AICc={-}7.99   BIC={-}6.12}

\NormalTok{monthly\_arima }\OtherTok{\textless{}{-}} \FunctionTok{auto.arima}\NormalTok{(monthly\_train)}

\CommentTok{\# Series: monthly\_train }
\CommentTok{\# ARIMA(2,0,2) with non{-}zero mean }
\CommentTok{\# }
\CommentTok{\# Coefficients:}
\CommentTok{\#          ar1      ar2      ma1     ma2    mean}
\CommentTok{\#       1.6522  {-}0.9335  {-}1.1833  0.4651  3.2705}
\CommentTok{\# s.e.  0.0423   0.0444   0.1831  0.1728  0.0830}
\CommentTok{\# }
\CommentTok{\# sigma\^{}2 = 2.048:  log likelihood = {-}526.49}
\CommentTok{\# AIC=1064.98   AICc=1065.27   BIC=1087.14}


\CommentTok{\# Fit ETS models}
\NormalTok{yearly\_ets }\OtherTok{\textless{}{-}} \FunctionTok{ets}\NormalTok{(yearly\_train)}

\CommentTok{\# ETS(A,N,N) }
\CommentTok{\# }
\CommentTok{\# Call:}
\CommentTok{\# ets(y = yearly\_train)}
\CommentTok{\# }
\CommentTok{\#   Smoothing parameters:}
\CommentTok{\#     alpha = 0.9999 }
\CommentTok{\# }
\CommentTok{\#   Initial states:}
\CommentTok{\#     l = 8.2942 }
\CommentTok{\# }
\CommentTok{\#   sigma:  0.2218}
\CommentTok{\# }
\CommentTok{\#      AIC     AICc      BIC }
\CommentTok{\# 54.74317 55.23297 60.65405 }

\NormalTok{monthly\_ets }\OtherTok{\textless{}{-}} \FunctionTok{ets}\NormalTok{(monthly\_train)}

\CommentTok{\# ETS(A,N,N) }
\CommentTok{\# }
\CommentTok{\# Call:}
\CommentTok{\# ets(y = monthly\_train)}
\CommentTok{\# }
\CommentTok{\#   Smoothing parameters:}
\CommentTok{\#     alpha = 0.9999 }
\CommentTok{\# }
\CommentTok{\#   Initial states:}
\CommentTok{\#     l = 3.2657 }
\CommentTok{\# }
\CommentTok{\#   sigma:  1.8451}
\CommentTok{\# }
\CommentTok{\#      AIC     AICc      BIC }
\CommentTok{\# 2058.871 2058.953 2069.952 }
\end{Highlighting}
\end{Shaded}

\subsubsection{Step 5, model performance residual
diagnostics}\label{step-5-model-performance-residual-diagnostics}

\begin{Shaded}
\begin{Highlighting}[]
\FunctionTok{checkresiduals}\NormalTok{(yearly\_arima)}
\end{Highlighting}
\end{Shaded}

\includegraphics{Lab1-ARIMA-2025-Team3_files/figure-latex/unnamed-chunk-6-1.pdf}

\begin{verbatim}
## 
##  Ljung-Box test
## 
## data:  Residuals from ARIMA(0,1,0)
## Q* = 5.659, df = 10, p-value = 0.8431
## 
## Model df: 0.   Total lags used: 10
\end{verbatim}

\begin{Shaded}
\begin{Highlighting}[]
\CommentTok{\# ACF of residuals shows no significant autocorrelation and looks like a normal distribution, indicating a very good fit}
\FunctionTok{checkresiduals}\NormalTok{(monthly\_arima)}
\end{Highlighting}
\end{Shaded}

\includegraphics{Lab1-ARIMA-2025-Team3_files/figure-latex/unnamed-chunk-6-2.pdf}

\begin{verbatim}
## 
##  Ljung-Box test
## 
## data:  Residuals from ARIMA(3,0,4)(0,1,1)[12]
## Q* = 30.958, df = 16, p-value = 0.01362
## 
## Model df: 8.   Total lags used: 24
\end{verbatim}

\begin{Shaded}
\begin{Highlighting}[]
\CommentTok{\# ACF of residuals have two autocorrelation spikes but could be random weird thing and looks like a normal distribution, indicating a relative good fit}
\FunctionTok{checkresiduals}\NormalTok{(yearly\_ets)}
\end{Highlighting}
\end{Shaded}

\includegraphics{Lab1-ARIMA-2025-Team3_files/figure-latex/unnamed-chunk-6-3.pdf}

\begin{verbatim}
## 
##  Ljung-Box test
## 
## data:  Residuals from ETS(A,N,N)
## Q* = 5.6202, df = 10, p-value = 0.8461
## 
## Model df: 0.   Total lags used: 10
\end{verbatim}

\begin{Shaded}
\begin{Highlighting}[]
\CommentTok{\# ACF of residuals shows no significant autocorrelation and looks like a normal distribution, indicating a good fit}
\FunctionTok{checkresiduals}\NormalTok{(monthly\_ets)}
\end{Highlighting}
\end{Shaded}

\includegraphics{Lab1-ARIMA-2025-Team3_files/figure-latex/unnamed-chunk-6-4.pdf}

\begin{verbatim}
## 
##  Ljung-Box test
## 
## data:  Residuals from ETS(A,N,A)
## Q* = 253.79, df = 24, p-value < 2.2e-16
## 
## Model df: 0.   Total lags used: 24
\end{verbatim}

\begin{Shaded}
\begin{Highlighting}[]
\CommentTok{\# ACF of residuals shows several significant autocorrelations though looks like a normal distribution, indicating a poor fit}

\DocumentationTok{\#\# In general, the ARIMA models seem to fit better than the ETS models for both yearly and monthly data.}
\end{Highlighting}
\end{Shaded}

\subsubsection{Step 6, forecast}\label{step-6-forecast}

\begin{Shaded}
\begin{Highlighting}[]
\CommentTok{\# Generate forecasts for the test period}
\NormalTok{yearly\_arima\_forecast }\OtherTok{\textless{}{-}} \FunctionTok{forecast}\NormalTok{(yearly\_arima, }\AttributeTok{h =} \FunctionTok{length}\NormalTok{(yearly\_test))}
\NormalTok{yearly\_ets\_forecast }\OtherTok{\textless{}{-}} \FunctionTok{forecast}\NormalTok{(yearly\_ets, }\AttributeTok{h =} \FunctionTok{length}\NormalTok{(yearly\_test))}

\NormalTok{monthly\_arima\_forecast }\OtherTok{\textless{}{-}} \FunctionTok{forecast}\NormalTok{(monthly\_arima, }\AttributeTok{h =} \FunctionTok{length}\NormalTok{(monthly\_test))}
\NormalTok{monthly\_ets\_forecast }\OtherTok{\textless{}{-}} \FunctionTok{forecast}\NormalTok{(monthly\_ets, }\AttributeTok{h =} \FunctionTok{length}\NormalTok{(monthly\_test))}

\CommentTok{\# Plot forecasts}
\NormalTok{plot1}\OtherTok{\textless{}{-}}\FunctionTok{autoplot}\NormalTok{(yearly\_arima\_forecast)}
\NormalTok{plot2}\OtherTok{\textless{}{-}}\FunctionTok{autoplot}\NormalTok{(yearly\_ets\_forecast)}

\NormalTok{plot3}\OtherTok{\textless{}{-}}\FunctionTok{autoplot}\NormalTok{(monthly\_arima\_forecast)}
\NormalTok{plot4}\OtherTok{\textless{}{-}}\FunctionTok{autoplot}\NormalTok{(monthly\_ets\_forecast)}

\CommentTok{\# Put into a 4 plot panel }
\FunctionTok{library}\NormalTok{(gridExtra)}
\FunctionTok{library}\NormalTok{(grid)}
\FunctionTok{library}\NormalTok{(ggplot2)}

\FunctionTok{grid.arrange}\NormalTok{(plot1, plot2, plot3, plot4, }\AttributeTok{ncol =} \DecValTok{2}\NormalTok{)}
\end{Highlighting}
\end{Shaded}

\includegraphics{Lab1-ARIMA-2025-Team3_files/figure-latex/unnamed-chunk-7-1.pdf}

\subsection{Accuracy metrics}\label{accuracy-metrics}

\begin{Shaded}
\begin{Highlighting}[]
\CommentTok{\# Evaluate forecast accuracy}
\NormalTok{yearly\_arima\_accuracy }\OtherTok{\textless{}{-}} \FunctionTok{accuracy}\NormalTok{(yearly\_arima\_forecast, yearly\_test)}

\CommentTok{\#                     ME      RMSE       MAE        MPE     MAPE      MASE      ACF1 Theil\textquotesingle{}s U}
\CommentTok{\# Training set {-}0.01046748 0.2175570 0.1732156 {-}0.1712576 2.225795 0.9820193 0.0292464        NA}
\CommentTok{\# Test set     {-}0.07686666 0.2753066 0.2363580 {-}1.1236585 3.107135 1.3399954 0.5697271  1.139735}

\NormalTok{yearly\_ets\_accuracy }\OtherTok{\textless{}{-}} \FunctionTok{accuracy}\NormalTok{(yearly\_ets\_forecast, yearly\_test)}

\CommentTok{\#                     ME      RMSE       MAE        MPE     MAPE      MASE       ACF1}
\CommentTok{\# Training set {-}0.01062423 0.2175548 0.1730609 {-}0.1731526 2.223930 0.9811418 0.02860653}
\CommentTok{\# Test set     {-}0.07683952 0.2752990 0.2363503 {-}1.1233035 3.107024 1.3399514 0.56972710}
             
\NormalTok{monthly\_arima\_accuracy }\OtherTok{\textless{}{-}} \FunctionTok{accuracy}\NormalTok{(monthly\_arima\_forecast, monthly\_test)}

\CommentTok{\#                     ME     RMSE       MAE  MPE MAPE      MASE         ACF1 Theil\textquotesingle{}s U}
\CommentTok{\# Training set {-}0.0620450 1.314323 0.9435214 {-}Inf  Inf 0.8521606 {-}0.001970859        NA}
\CommentTok{\# Test set     {-}0.9640865 3.266402 2.5277314 {-}Inf  Inf 2.2829721  0.737652808         0}

\NormalTok{monthly\_ets\_accuracy }\OtherTok{\textless{}{-}} \FunctionTok{accuracy}\NormalTok{(monthly\_ets\_forecast, monthly\_test)}

\CommentTok{\#                       ME     RMSE      MAE  MPE MAPE     MASE      ACF1 Theil\textquotesingle{}s U}
\CommentTok{\# Training set  0.0008512003 1.653080 1.088517 {-}Inf  Inf 0.983116 0.3694314        NA}
\CommentTok{\# Test set     {-}1.8938374286 3.914736 2.998623 {-}Inf  Inf 2.708267 0.7424152         0}
\end{Highlighting}
\end{Shaded}

\subsection{Results and Discussion}\label{results-and-discussion}

The accuracy metrics show that for yearly data, ARIMA and ETS models
performed nearly the same with a MAPE value of about 3.1\%, indicating
reasonably accurate forecasts. For the monthly data, the ARIMA model
outperformed ETS with lower RMSE (3.27 vs 3.91) and MAE (2.53 vs 3.00)

Visually, the plots suggest better monthly performance but the numerical
metrics indicate higher absolute errors for monthly forecasts compared
to yearly ones. It would be interesting to compare how each of these
predictions preform on test data from the other dataset (i.e.~do
predictions made from monthly data preform better for yearly data).
Similarly, it would be interesting to compare the results here to that
of the data from California, which is made complicated by the different
range of time series values. We questioned the legitimacy of our methods
for fitting an ARIMA to the yearly data, as the model suggests that the
predicted catch for the next years is equal to the last year's catch.
This held true with some back-of-the envelope calculations using the
same methods on the California and Oregon data. Specifically, we
question the fact that the forecast plot visually shows no variation
that would typically be associated with a random walk, rather
forecasting a straight line with no variation.

\#\#Team Member Contributions

BM, EC, and TT all contributed code, question design, and analysis to
the final work presented here. EC wrote details about the initial plan
and questions we addressed, and how the plan was carried out, as well as
added details to the discussion section. TT added details about checking
the residuals and the ACF/PACF plots, and wrote details about the
results and discussion. BM wrote and organized final code and added
details about Accuracy results.

\end{document}
